\usepackage[margin=1in]{geometry}
\usepackage{amsmath,amssymb,amsthm,mathtools}
\usepackage{booktabs}

\newtheorem{theorem}{Theorem}[section]
\newtheorem{lemma}[theorem]{Lemma}
\newtheorem{proposition}[theorem]{Proposition}
\newtheorem{corollary}[theorem]{Corollary}
\theoremstyle{definition}
\newtheorem{definition}[theorem]{Definition}
\newtheorem{certificate}[theorem]{Certificate}
\theoremstyle{remark}
\newtheorem{remark}[theorem]{Remark}

\newcommand{\J}{\mathcal J}
\newcommand{\Gm}{\mathcal G}
\newcommand{\Lm}{\mathcal L}
\newcommand{\Em}{\mathcal E}
\newcommand{\C}{\mathbb C}
\newcommand{\R}{\mathbb R}
\newcommand{\N}{\mathbb N}
\newcommand{\dd}{\,\mathrm d}
\newcommand{\boxconv}{\boxtimes_d}
\newcommand{\norm}[1]{\left\lVert #1\right\rVert}
\newcommand{\abs}[1]{\left\lvert #1\right\rvert}
\newcommand{\sech}{\operatorname{sech}}
\newcommand{\lmesh}{\operatorname{lmesh}}
\newcommand{\dist}{\operatorname{dist}}
\newcommand{\code}[1]{\texttt{#1}}

\allowdisplaybreaks
\numberwithin{equation}{section}
